\chapter{家用化学品卫生}
\setcounter{chapter}{9}
\chapterintro{
了解化妆品的定义、分类及对健康的不良影响；掌握化妆品对皮肤的损害类型；了解化妆品微生物污染的危害。
}

\section{化妆品概述}

\begin{defbox}
\textbf{化妆品（Cosmetic）}：以涂擦、喷洒或其他类似方法，散布于人体任何表面（皮肤、毛发、指甲、口唇等）、牙齿及口腔黏膜，以清洁、保护、美化、修饰及保持其处于良好状态为目的的产品。
\end{defbox}

\subsection{化妆品分类}

\begin{enumerate}[leftmargin=*]
    \item \imp{一般用途化妆品}：护肤类、益发类、美容修饰类、芳香类、口腔卫生用品
    \item \imp{特殊用途化妆品}：用于育发、染发、烫发、脱毛、除臭、祛斑和防晒的化妆品。常需加入某些活性物质以弥补体表局部缺陷，其中有些有一定副作用而在化妆品中被限制使用
\end{enumerate}

\section{化妆品对健康的不良影响}

\subsection{化妆品对皮肤的不良影响}

\begin{keybox}
\textbf{1. 刺激性接触性皮炎（ICD）}：\imp{最常见}的化妆品皮肤损害，因反复接触引起。无免疫记忆，首次接触即可发生。

\textbf{2. 变应性接触性皮炎}：接触化妆品数天后局部出现皮炎症状，属于迟发型变态反应，第二次接触后发生。

\textbf{3. 化妆品光感性皮炎}：化妆品中某些成分在日光（紫外线）作用下引起的光毒性或光变应性反应。包括光变应性接触性皮炎（PCD）和光毒性皮炎。

\textbf{4. 化妆品痤疮}：由矿物油制成的基质引起，油性皮肤者多见。诊断需重点了解化妆品接触史和施用部位。

\textbf{5. 化妆品皮肤色素异常}：化妆品成分导致的色素沉着或色素减退。
\end{keybox}

\subsection{化妆品微生物污染危害}

\begin{enumerate}[leftmargin=*]
    \item \imp{一次污染}：化妆品原料本身有微生物，或在生产过程中受污染
    \item \imp{二次污染}：化妆品启封后，使用或存放过程中受到污染，包括手部接触带入微生物或空气中微生物落入
\end{enumerate}

微生物污染可引起化妆品腐败变质，其在代谢过程中产生的毒素作为变应原或刺激原对施用部位产生致敏或刺激作用。此外，化妆品被致病菌污染可能引起局部甚至全身感染。

\subsection{化妆品所含化学物质的毒性作用}

特殊用途化妆品中常需加入某些活性物质（如染发剂中的对苯二胺、烫发剂中的巯基乙酸等），这些化学物质可能通过皮肤吸收产生全身毒性作用，如致突变、致癌、致畸等。

\section{复习题}

\begin{enumerate}[leftmargin=*, itemsep=8pt]
    \item \textbf{【名词解释】}化妆品；ICD
    \item \textbf{【简答题】}简述化妆品的分类。
    \item \textbf{【简答题】}简述化妆品对皮肤的主要不良影响类型。
    \item \textbf{【简答题】}简述化妆品微生物污染的一次污染与二次污染的区别。
\end{enumerate}

\subsection*{参考答案要点}

\begin{enumerate}[leftmargin=*, itemsep=5pt]
    \item \textbf{化妆品}：以涂擦、喷洒或其他类似方法散布于人体表面、牙齿及口腔黏膜，以清洁、保护、美化、修饰及保持良好状态为目的的产品。\textbf{ICD}：刺激性接触性皮炎，化妆品最常见的皮肤损害。
    \item 一般用途（护肤、益发、美容修饰、芳香、口腔卫生）和特殊用途（育发、染发、烫发、脱毛、除臭、祛斑、防晒）。
    \item 五种类型：刺激性接触性皮炎（ICD，最常见）、变应性接触性皮炎、化妆品光感性皮炎（PCD和光毒性皮炎）、化妆品痤疮、化妆品皮肤色素异常。
    \item 一次污染：原料本身有微生物或生产过程中受污染。二次污染：启封后使用或存放过程中受污染（手部接触或空气中微生物落入）。
\end{enumerate}

\newpage

% =====================================================================
%  CHAPTER 10: 公共场所卫生
% =====================================================================
\newpage